% =============================================================================
% R-20 Reviewer Response Materials
% 评价指标完整性补全：判别力（ROC-AUC/PR-AUC/MCC）+ 校准（Brier/ECE）
%                     + 候选排序能力（Spearman + Top-k enrichment）
%
% 用法：本文件为草稿与表格模板，数值字段以 \xx 占位（待 5fold 跑完后填入）。
%       可按需把对应片段粘贴进主稿对应小节。
% =============================================================================

% -----------------------------------------------------------------------------
% 1) 评价指标小节改写（替换原 p11 L21-L24 那段评价体系描述）
% -----------------------------------------------------------------------------
为了全面评估模型的分类性能、概率预测能力以及候选序列排序能力，本文采用
多维评价指标体系。分类性能方面，采用 Subset Accuracy、Example-level
Accuracy/Precision/Recall/F1 和 Label-level Accuracy/Precision/Recall/F1
衡量不同粒度下的预测准确性。进一步地，引入 ROC-AUC、PR-AUC 和 Matthews
Correlation Coefficient (MCC) 评价模型在不同决策阈值下的判别能力。此外，
通过 Brier score 和 Expected Calibration Error (ECE) 分析预测概率的可靠性。
针对模型在候选序列筛选中的应用场景，本文进一步计算 predicted cleavage ratio
与实验观测断裂比例之间的 Spearman correlation，并开展 candidate-level ranking
test 验证模型输出作为候选优先级评分的有效性。

% -----------------------------------------------------------------------------
% 2) 新增小节：候选序列排序分析（建议放在主结果之后）
% -----------------------------------------------------------------------------
\subsection{Candidate-level ranking analysis}
\label{sec:candidate_ranking}

为验证 DBond-GT 的输出是否可作为镜像肽候选序列的优先级评分，本文设计了
candidate-level ranking test。对于每条肽，模型输出所有候选键的断裂概率
$p_1, p_2, \ldots, p_{n}$，定义 peptide-level predicted cleavage ratio 为
键级概率的均值：
\begin{equation}
    R_{\text{pred}} = \frac{1}{n}\sum_{i=1}^{n} p_i.
    \label{eq:rpred}
\end{equation}
对应地，实验观测的 PBCLA cleavage ratio 为真实断裂键数与总键数之比：
\begin{equation}
    R_{\text{true}} = \frac{1}{n}\sum_{i=1}^{n} y_i.
    \label{eq:rtrue}
\end{equation}
$R_{\text{pred}}$ 越高表示模型认为该肽段越容易碎裂、越适合测序。我们将全部
测试肽按 $R_{\text{pred}}$ 降序排列，并计算 $R_{\text{pred}}$ 与 $R_{\text{true}}$
之间的 Spearman 秩相关 $\rho$，同时在 Top 10\%/20\%/50\% 候选子集上比较
观测断裂率与全集基线的差异，以评估排序的下游筛选价值（图~\ref{fig:ranking}）。

% -----------------------------------------------------------------------------
% 3) 三张结果表模板（5fold mean ± std）
% -----------------------------------------------------------------------------

% 表 1：判别力指标（阈值无关）
\begin{table}[t]
\centering
\caption{Discrimination metrics of DBond-GT on the 5-fold cross-validation
(mean $\pm$ std). AUC and PR-AUC measure threshold-free ranking quality;
MCC summarizes the confusion matrix in a single imbalance-robust score.}
\label{tab:discrimination}
\begin{tabular}{lccc}
\toprule
Model     & ROC-AUC         & PR-AUC          & MCC             \\
\midrule
DBond-GT  & \xx $\pm$ \xx   & \xx $\pm$ \xx   & \xx $\pm$ \xx   \\
\bottomrule
\end{tabular}
\end{table}

% 表 2：概率校准指标
\begin{table}[t]
\centering
\caption{Probability calibration of DBond-GT. Brier score is the mean squared
error of probabilistic predictions; ECE is the expected calibration error over
10 equal-width bins. Lower is better for both.}
\label{tab:calibration}
\begin{tabular}{lcc}
\toprule
Model     & Brier score     & ECE             \\
\midrule
DBond-GT  & \xx $\pm$ \xx   & \xx $\pm$ \xx   \\
\bottomrule
\end{tabular}
\end{table}

% 表 3：候选排序能力（R-20 核心）
\begin{table}[t]
\centering
\caption{Candidate-level ranking capability of DBond-GT. Spearman $\rho$ is
computed between predicted and observed cleavage ratios across all test
peptides; Top-$k$ enrichment reports the observed cleavage ratio of the
top-ranked subset relative to the all-peptide baseline. Higher $\rho$ and
larger enrichment indicate the predicted cleavage ratio is an effective
candidate ranking score.}
\label{tab:ranking}
\begin{tabular}{lcccc}
\toprule
Model     & Spearman $\rho$   & Top 10\%       & Top 20\%       & Top 50\%       \\
\midrule
DBond-GT  & \xx (p=\xx)       & \xx            & \xx            & \xx            \\
\midrule
All (baseline) & ---          & \xx            & \xx            & \xx            \\
\bottomrule
\end{tabular}
\end{table}

% 图引用
\begin{figure}[t]
\centering
% \includegraphics[width=\textwidth]{figures/panel_reliability_diagram.pdf}
% \includegraphics[width=\textwidth]{figures/panel_predicted_vs_observed.pdf}
% \includegraphics[width=\textwidth]{figures/panel_topk_enrichment.pdf}
\caption{\textbf{Probability calibration and candidate-level ranking of
DBond-GT.} (a) Reliability diagram: predicted bond-cleavage probabilities
closely match observed cleavage frequencies (ECE~$=$~\xx, Brier~$=$~\xx).
(b) Scatter of peptide-level predicted vs.\ observed cleavage ratios
(Spearman $\rho=$~\xx, $p<$~\xx). (c) Top-$k$ enrichment curve: peptides
ranked by $R_{\text{pred}}$ show monotonically higher observed cleavage
ratios than the all-peptide baseline, confirming the predicted cleavage
ratio is an effective candidate ranking score.}
\label{fig:ranking}
\end{figure}

% -----------------------------------------------------------------------------
% 4) Point-by-point 审稿回应英文段落（response letter 用）
% -----------------------------------------------------------------------------
% Reviewer R-20 [Moderate] — Evaluation metric completeness (p11 L21-L24)
%
\textbf{Response:} We thank the reviewer for highlighting the importance of
ranking and calibration for downstream sequence screening. We have expanded
the evaluation framework along three complementary dimensions:

\begin{enumerate}
    \item \textbf{Threshold-free discrimination.} In addition to fixed-threshold
    classification metrics, we now report ROC-AUC, PR-AUC, and Matthews
    Correlation Coefficient (MCC) to characterize the model's ability to
    distinguish cleavage from non-cleavage bonds across all operating points
    (Table~\ref{tab:discrimination}).

    \item \textbf{Probability calibration.} We added Brier score and Expected
    Calibration Error (ECE) to assess whether the predicted bond-cleavage
    probabilities are trustworthy in absolute terms (Table~\ref{tab:calibration}
    and Fig.~\ref{fig:ranking}a).

    \item \textbf{Candidate-level ranking.} To directly address whether the
    predicted cleavage ratio can serve as a candidate ranking score, we
    aggregate bond-level probabilities into a peptide-level predicted cleavage
    ratio $R_{\text{pred}}$ (Eq.~\ref{eq:rpred}) and compare it with the
    PBCLA-observed ratio $R_{\text{true}}$ (Eq.~\ref{eq:rtrue}). The Spearman
    correlation reaches $\rho =$~\xx ($p <$~\xx), and Top-$k$ enrichment
    analysis confirms that peptides ranked highest by $R_{\text{pred}}$ exhibit
    systematically higher observed cleavage ratios than the baseline
    (Table~\ref{tab:ranking}, Fig.~\ref{fig:ranking}b--c). These results
    demonstrate that DBond-GT outputs are effective candidate ranking scores
    for mirror-image peptide screening.
\end{enumerate}
